% =====================================================================
% Lösungsteil zu Kapitel 2 (Analytische Geometrie) – Lösungen zu Ü 2.1 bis Ü 2.22
% Steht am Ende des Buches hinter dem Anhang.
% Alle Ergebnisse wurden mit SymPy nachgerechnet.
% =====================================================================
\Loesungskapitel{Lösungen zur Analytischen Geometrie}

\Uebungsabschnitt{Vektoren und Vektorraum}

\begin{loesung}{2.1}{Vektorketten im Spat}
\textit{Grundlage:} Im Spat sind gegenüberliegende Kanten parallel und gleich lang. Es gilt also $\vv{BC}=\vv{AD}=\vv{EH}=\vv{FG}=\vv{b}$, $\vv{DC}=\vv{EF}=\vv{HG}=\vv{a}$ und $\vv{BF}=\vv{CG}=\vv{DH}=\vv{c}$. Ein Verbindungsvektor wird als Vektorkette entlang der Kanten geschrieben.
\begin{enumerate}[label=\alph*)]
  \item
  \begin{align*}
    \vv{AG} &= \vv{AB}+\vv{BC}+\vv{CG} = \vv{a}+\vv{b}+\vv{c}\\
    \vv{BH} &= \vv{BA}+\vv{AD}+\vv{DH} = -\vv{a}+\vv{b}+\vv{c}\\
    \vv{CE} &= \vv{CD}+\vv{DA}+\vv{AE} = -\vv{a}-\vv{b}+\vv{c}
  \end{align*}
  \item $M$ halbiert $\overline{FH}$, also $\vv{FM} = \frac{1}{2}\vv{FH}$ mit $\vv{FH} = \vv{FE}+\vv{EH} = -\vv{a}+\vv{b}$.
  \[\vv{AM} = \vv{AB}+\vv{BF}+\vv{FM} = \vv{a}+\vv{c}+\dfrac{1}{2}(-\vv{a}+\vv{b}) = \dfrac{1}{2}\vv{a}+\dfrac{1}{2}\vv{b}+\vv{c}\]
  \item \textit{Idee:} Es reicht zu zeigen, dass die Mittelpunkte $M_1$ von $\overline{AG}$ und $M_2$ von $\overline{BH}$ zusammenfallen.
  \begin{align*}
    \vv{AM_1} &= \dfrac{1}{2}\vv{AG} = \dfrac{1}{2}(\vv{a}+\vv{b}+\vv{c})\\
    \vv{AM_2} &= \vv{AB}+\dfrac{1}{2}\vv{BH} = \vv{a}+\dfrac{1}{2}(-\vv{a}+\vv{b}+\vv{c}) = \dfrac{1}{2}(\vv{a}+\vv{b}+\vv{c})
  \end{align*}
  Wegen $\vv{AM_1}=\vv{AM_2}$ gilt $M_1=M_2$. Der Mittelpunkt beider Raumdiagonalen ist derselbe Punkt, die Diagonalen schneiden sich dort und halbieren sich gegenseitig.
\end{enumerate}
\end{loesung}

\begin{loesung}{2.2}{Orts- und Verbindungsvektoren}
\begin{enumerate}[label=\alph*)]
  \item Verbindungsvektor \glqq Spitze minus Fuß\grqq, Länge $|\vv{v}| = \sqrt{v_1^2+v_2^2+v_3^2}$:
  \begin{align*}
    \vv{AB} &= \begin{pmatrix} 2\\ 2\\ -1\end{pmatrix}, & |\vv{AB}| &= \sqrt{4+4+1} = 3\\
    \vv{BC} &= \begin{pmatrix} -2\\ 2\\ 1\end{pmatrix}, & |\vv{BC}| &= \sqrt{4+4+1} = 3\\
    \vv{AC} &= \begin{pmatrix} 0\\ 4\\ 0\end{pmatrix}, & |\vv{AC}| &= 4
  \end{align*}
  \item Wegen $|\vv{AB}| = |\vv{BC}| = 3$ sind zwei Seiten gleich lang, das Dreieck ist gleichschenklig. Die Schenkel sind $\overline{AB}$ und $\overline{BC}$, die Basis ist $\overline{AC}$ mit der Länge $4$. Die Spitze ist $B$.
  \item Schwerpunkt des Dreiecks (Merkekasten): $\vv{S} = \frac{1}{3}(\vv{A}+\vv{B}+\vv{C})$.
  \[\vv{S} = \dfrac{1}{3}\begin{pmatrix} 2+4+2\\ 0+2+4\\ 1+0+1\end{pmatrix} = \begin{pmatrix} \frac{8}{3}\\ 2\\ \frac{2}{3}\end{pmatrix},\qquad S\left(\tfrac{8}{3}\big|2\big|\tfrac{2}{3}\right)\]
  \item Im Parallelogramm $ABCD$ gilt $\vv{AD} = \vv{BC}$.
  \[\vv{D} = \vv{A}+\vv{BC} = \begin{pmatrix} 2\\ 0\\ 1\end{pmatrix}+\begin{pmatrix} -2\\ 2\\ 1\end{pmatrix} = \begin{pmatrix} 0\\ 2\\ 2\end{pmatrix},\qquad D(0|2|2)\]
  Im Parallelogramm sind gegenüberliegende Seiten gleich lang. Aus $|\vv{AB}| = |\vv{BC}| = 3$ folgt, dass alle vier Seiten die Länge $3$ haben. Damit ist $ABCD$ eine Raute.
\end{enumerate}
\end{loesung}

\Uebungsabschnitt{Skalar-, Vektor- und Spatprodukt}

\begin{loesung}{2.3}{Skalarprodukt und Winkel}
\begin{enumerate}[label=\alph*)]
  \item Der Winkel $\beta$ wird von den Vektoren $\vv{BA}$ und $\vv{BC}$ eingeschlossen (beide gehen von $B$ aus). Winkel zwischen Vektoren: $\cos{\beta} = \dfrac{\vv{BA}\circ\vv{BC}}{|\vv{BA}|\cdot|\vv{BC}|}$.
  \[\vv{BA} = \begin{pmatrix} -2\\ -2\\ 1\end{pmatrix},\quad \vv{BC} = \begin{pmatrix} -2\\ 2\\ 1\end{pmatrix},\quad \vv{BA}\circ\vv{BC} = 4-4+1 = 1\]
  \[\cos{\beta} = \dfrac{1}{3\cdot 3} = \dfrac{1}{9}\quad\Longrightarrow\quad\beta\approx 83{,}62^{\circ}\]
  \item Nach Ü~2.2\,d) ist $ABCD$ eine Raute. Eine Raute ist genau dann ein Quadrat, wenn sie einen rechten Winkel hat. Orthogonalitätskriterium: $\vv{u}\perp\vv{v}\Longleftrightarrow\vv{u}\circ\vv{v}=0$.
  \[\vv{AB}\circ\vv{BC} = 2\cdot(-2)+2\cdot 2+(-1)\cdot 1 = -1\neq 0\]
  Bei $B$ liegt kein rechter Winkel vor, $ABCD$ ist also \textbf{kein} Quadrat. (Das passt zu a): $\beta\approx 83{,}62^{\circ}$.)
  \item Orthogonalitätskriterium:
  \[\begin{pmatrix} 1\\ t\\ 2\end{pmatrix}\circ\begin{pmatrix} 3\\ -1\\ t\end{pmatrix} = 3-t+2t = 3+t = 0\quad\Longleftrightarrow\quad t = -3\]
\end{enumerate}
\end{loesung}

\begin{loesung}{2.4}{Vektorprodukt und Spatprodukt}
\begin{enumerate}[label=\alph*)]
  \item $\vv{AB} = \begin{pmatrix} -1\\ 2\\ 0\end{pmatrix}$, $\vv{AC} = \begin{pmatrix} -1\\ 0\\ 3\end{pmatrix}$. Berechnung des Vektorprodukts:
  \[\vv{n} = \vv{AB}\times\vv{AC} = \begin{pmatrix} 2\cdot 3-0\cdot 0\\ 0\cdot(-1)-(-1)\cdot 3\\ (-1)\cdot 0-2\cdot(-1)\end{pmatrix} = \begin{pmatrix} 6\\ 3\\ 2\end{pmatrix}\]
  \textit{Nachweis der Orthogonalität} mit dem Skalarprodukt:
  \[\vv{n}\circ\vv{AB} = -6+6+0 = 0,\qquad \vv{n}\circ\vv{AC} = -6+0+6 = 0\]
  \item Der Betrag des Vektorprodukts ist der Flächeninhalt des aufgespannten Parallelogramms, das Dreieck hat die Hälfte davon.
  \[|\vv{n}| = \sqrt{36+9+4} = 7,\qquad A_{ABC} = \dfrac{1}{2}\cdot 7 = 3{,}5\ \text{FE}\]
  \item Die Pyramide $OABC$ wird von $\vv{OA}$, $\vv{OB}$ und $\vv{OC}$ aufgespannt. Definition Spatprodukt: $V_{\text{Pyramide}} = \frac{1}{6}\left|(\vv{a}\times\vv{b})\circ\vv{c}\right|$.
  \[\vv{OA}\times\vv{OB} = \begin{pmatrix} 1\\ 0\\ 0\end{pmatrix}\times\begin{pmatrix} 0\\ 2\\ 0\end{pmatrix} = \begin{pmatrix} 0\\ 0\\ 2\end{pmatrix},\qquad \begin{pmatrix} 0\\ 0\\ 2\end{pmatrix}\circ\begin{pmatrix} 0\\ 0\\ 3\end{pmatrix} = 6\]
  \[V = \dfrac{1}{6}\cdot|6| = 1\ \text{VE}\]
  \textit{Kontrolle:} Grundfläche $OAB$ mit $\frac{1}{2}\cdot 1\cdot 2 = 1$, Höhe $3$, also $V = \frac{1}{3}\cdot 1\cdot 3 = 1$.
\end{enumerate}
\end{loesung}

\Uebungsabschnitt{Lineare Abhängigkeit und Unabhängigkeit}

\begin{loesung}{2.5}{Lineare Abhängigkeit}
\textit{Grundlage:} Drei Vektoren sind linear abhängig, wenn sich einer als Linearkombination der anderen darstellen lässt.
\begin{enumerate}[label=\alph*)]
  \item \[\begin{pmatrix} 1\\ 2\\ 3\end{pmatrix}+\begin{pmatrix} 2\\ 3\\ 1\end{pmatrix} = \begin{pmatrix} 3\\ 5\\ 4\end{pmatrix}\]
  Der dritte Vektor ist die Summe der beiden ersten, die drei Vektoren sind also linear abhängig.
  \item Die Vektoren $\begin{pmatrix} 0\\ 1\\ 2\end{pmatrix}$ und $\begin{pmatrix} 1\\ 1\\ 1\end{pmatrix}$ sind nicht parallel (die erste Koordinate ist $0$ bzw. $1$, die zweite jeweils $1$). Die drei Vektoren sind also genau dann linear abhängig, wenn sich der erste Vektor aus diesen beiden kombinieren lässt:
  \[r\cdot\begin{pmatrix} 0\\ 1\\ 2\end{pmatrix}+s\cdot\begin{pmatrix} 1\\ 1\\ 1\end{pmatrix} = \begin{pmatrix} 1\\ 0\\ a\end{pmatrix}\qquad\begin{aligned} \text{I: }& s = 1\\ \text{II: }& r+s = 0\\ \text{III: }& 2r+s = a\end{aligned}\]
  Aus I und II folgt $s=1$, $r=-1$. In III: $a = -2+1 = -1$. Für $a=-1$ sind die Vektoren linear abhängig.
  \item Drei linear abhängige Vektoren liegen, als Pfeile von einem gemeinsamen Punkt aus abgetragen, in einer gemeinsamen Ebene (sie sind komplanar). Sie spannen dann keinen Spat auf, und ihr Spatprodukt ist $0$. Drei linear unabhängige Vektoren spannen dagegen den ganzen Raum auf.
\end{enumerate}
\end{loesung}

\Uebungsabschnitt{Geraden im Raum}

\begin{loesung}{2.6}{Spurpunkte, Lagebeziehung und Schnittwinkel von Geraden}
Ortsvektor eines allgemeinen Punktes von $g$: $\vv{X} = \begin{pmatrix} 1+\lambda\\ 2-\lambda\\ 3+2\lambda\end{pmatrix}$.
\begin{enumerate}[label=\alph*)]
  \item \textit{Verfahren:} Merkekasten \glqq Berechnung der Spurpunkte\grqq, d.\,h. jeweils eine Koordinate null setzen.
  \begin{itemize}
    \item $x_1x_2$-Ebene ($x_3=0$): $3+2\lambda = 0$, also $\lambda = -\frac{3}{2}$ und $S_{12}\left(-\frac{1}{2}\big|\frac{7}{2}\big|0\right)$.
    \item $x_1x_3$-Ebene ($x_2=0$): $2-\lambda = 0$, also $\lambda = 2$ und $S_{13}(3|0|7)$.
    \item $x_2x_3$-Ebene ($x_1=0$): $1+\lambda = 0$, also $\lambda = -1$ und $S_{23}(0|3|1)$.
  \end{itemize}
  \item \textit{Vorgehen:} Zuerst die Richtungsvektoren auf Parallelität prüfen, dann ein Gleichungssystem lösen.
  \begin{itemize}
    \item \textbf{$g$ und $h_1$:} $\begin{pmatrix} -2\\ 2\\ -4\end{pmatrix} = -2\cdot\begin{pmatrix} 1\\ -1\\ 2\end{pmatrix}$, die Richtungsvektoren sind also parallel. Liegt der Aufpunkt $O(0|0|0)$ von $h_1$ auf $g$? Aus $1+\lambda = 0$ folgt $\lambda = -1$, dann ist aber $x_2 = 2-(-1) = 3\neq 0$. $O$ liegt nicht auf $g$, die Geraden sind \textbf{echt parallel}.
    \item \textbf{$g$ und $h_2$:} Die Richtungsvektoren sind nicht parallel (erste Koordinate $1$ bzw. $0$). Gleichsetzen:
    \[\begin{aligned} \text{I: }& 1+\lambda = 2\\ \text{II: }& 2-\lambda = 3+\mu\\ \text{III: }& 3+2\lambda = 7+\mu\end{aligned}\qquad\Longrightarrow\qquad \lambda = 1,\ \mu = -2\]
    Probe in III: $3+2 = 5 = 7-2$, wahr. Die Geraden \textbf{schneiden sich} in $S(2|1|5)$ ($\lambda=1$ in $g$ eingesetzt).
    \item \textbf{$g$ und $h_3$:} Die Richtungsvektoren sind nicht parallel ($x_3$-Koordinate $2$ bzw. $0$, aber nicht beide null). Gleichsetzen:
    \[\begin{aligned} \text{I: }& 1+\lambda = \mu\\ \text{II: }& 2-\lambda = 1+\mu\\ \text{III: }& 3+2\lambda = 0\end{aligned}\]
    Aus III: $\lambda = -\frac{3}{2}$, aus I: $\mu = -\frac{1}{2}$. Probe in II: $2+\frac{3}{2} = \frac{7}{2}$, aber $1-\frac{1}{2} = \frac{1}{2}$, Widerspruch. Es gibt keinen Schnittpunkt, und die Geraden sind nicht parallel. Sie sind also \textbf{windschief}.
  \end{itemize}
  \item Schnittwinkel von Geraden (nicht stumpfer Winkel, daher Betrag im Zähler):
  \[\cos{\varphi} = \dfrac{\left|\begin{pmatrix} 1\\ -1\\ 2\end{pmatrix}\circ\begin{pmatrix} 0\\ 1\\ 1\end{pmatrix}\right|}{\sqrt{6}\cdot\sqrt{2}} = \dfrac{|0-1+2|}{\sqrt{12}} = \dfrac{1}{\sqrt{12}}\quad\Longrightarrow\quad\varphi\approx 73{,}22^{\circ}\]
\end{enumerate}
\end{loesung}

\Uebungsabschnitt{Ebenen im Raum}

\begin{loesung}{2.7}{Darstellungsformen von Ebenen}
\begin{enumerate}[label=\alph*)]
  \item Aufpunkt $A$, Spannvektoren $\vv{AB}$ und $\vv{AC}$ (nicht parallel):
  \[E:\ \vv{X} = \begin{pmatrix} 2\\ 0\\ 0\end{pmatrix}+\lambda\cdot\begin{pmatrix} -2\\ 4\\ 0\end{pmatrix}+\mu\cdot\begin{pmatrix} -2\\ 0\\ 4\end{pmatrix},\qquad\lambda,\mu\in\mathds{R}\]
  \item Ein Normalenvektor steht senkrecht auf beiden Spannvektoren, er ergibt sich also aus dem Vektorprodukt:
  \[\vv{AB}\times\vv{AC} = \begin{pmatrix} 4\cdot 4-0\cdot 0\\ 0\cdot(-2)-(-2)\cdot 4\\ (-2)\cdot 0-4\cdot(-2)\end{pmatrix} = \begin{pmatrix} 16\\ 8\\ 8\end{pmatrix} = 8\cdot\begin{pmatrix} 2\\ 1\\ 1\end{pmatrix},\qquad \vv{n} = \begin{pmatrix} 2\\ 1\\ 1\end{pmatrix}\]
  Koordinatenform: $2x_1+x_2+x_3-c = 0$. Punktprobe mit $A$: $4-c=0$, also $c=4$ und $E: 2x_1+x_2+x_3-4=0$.\\
  \textit{Kontrolle mit $B$ und $C$:} $0+4+0-4 = 0$ und $0+0+4-4 = 0$.
  \item $2\cdot 1+1+1-4 = 0$, die Gleichung ist erfüllt, also $P\in E$.
  \item \textit{Vorgehen:} Drei Punkte von $F$ suchen, die nicht auf einer Geraden liegen, oder einen Aufpunkt und zwei nicht parallele Vektoren, die senkrecht zu $\vv{n}_F = \begin{pmatrix} 1\\ -2\\ 2\end{pmatrix}$ sind.
  \begin{itemize}
    \item Aufpunkt: $x_2=x_3=0$ liefert $x_1=6$, also $Q(6|0|0)$.
    \item $\vv{u} = \begin{pmatrix} 2\\ 1\\ 0\end{pmatrix}$: $\vv{n}_F\circ\vv{u} = 2-2+0 = 0$, und $\vv{v} = \begin{pmatrix} 0\\ 1\\ 1\end{pmatrix}$: $\vv{n}_F\circ\vv{v} = 0-2+2 = 0$. Die Vektoren $\vv{u}$ und $\vv{v}$ sind nicht parallel.
  \end{itemize}
  \[F:\ \vv{X} = \begin{pmatrix} 6\\ 0\\ 0\end{pmatrix}+\lambda\cdot\begin{pmatrix} 2\\ 1\\ 0\end{pmatrix}+\mu\cdot\begin{pmatrix} 0\\ 1\\ 1\end{pmatrix},\qquad\lambda,\mu\in\mathds{R}\]
  (Andere Aufpunkte und Spannvektoren sind ebenfalls richtig.)
\end{enumerate}
\end{loesung}

\begin{loesung}{2.8}{Lagebeziehungen und Winkel}
Normalenvektoren: $\vv{n}_E = \begin{pmatrix} 2\\ 1\\ 1\end{pmatrix}$, $\vv{n}_F = \begin{pmatrix} 1\\ -2\\ 2\end{pmatrix}$.
\begin{enumerate}[label=\alph*)]
  \item \textit{Kriterium:} $g$ ist parallel zu $E$, wenn der Richtungsvektor von $g$ senkrecht auf $\vv{n}_E$ steht. $g$ ist echt parallel, wenn zusätzlich der Aufpunkt nicht in $E$ liegt.
  \[\vv{n}_E\circ\begin{pmatrix} 1\\ -1\\ -1\end{pmatrix} = 2-1-1 = 0\]
  Punktprobe mit dem Aufpunkt $(1|1|4)$: $2+1+4-4 = 3\neq 0$. Der Aufpunkt liegt nicht in $E$, $g$ verläuft also echt parallel zu $E$.
  \item Allgemeinen Punkt von $k$, $(\lambda|\lambda|1+\lambda)$, in $E$ einsetzen:
  \[2\lambda+\lambda+(1+\lambda)-4 = 0\quad\Longleftrightarrow\quad 4\lambda = 3\quad\Longleftrightarrow\quad\lambda = \dfrac{3}{4}\]
  \[S\left(\tfrac{3}{4}\big|\tfrac{3}{4}\big|\tfrac{7}{4}\right)\]
  Schnittwinkel Gerade--Ebene: $\sin{\varphi} = \dfrac{|\vv{u}\circ\vv{n}_E|}{|\vv{u}|\cdot|\vv{n}_E|}$, denn der Winkel zum Normalenvektor ergänzt $\varphi$ zu $90^{\circ}$.
  \[\sin{\varphi} = \dfrac{|2+1+1|}{\sqrt{3}\cdot\sqrt{6}} = \dfrac{4}{\sqrt{18}}\quad\Longrightarrow\quad\varphi\approx 70{,}53^{\circ}\]
  \item $\vv{n}_E$ und $\vv{n}_F$ sind nicht parallel ($2:1 \neq 1:(-2)$). Die Ebenen sind also nicht parallel und schneiden sich in einer Geraden.\\
  \textit{Richtungsvektor} der Schnittgeraden: senkrecht zu beiden Normalenvektoren.
  \[\vv{n}_E\times\vv{n}_F = \begin{pmatrix} 1\cdot 2-1\cdot(-2)\\ 1\cdot 1-2\cdot 2\\ 2\cdot(-2)-1\cdot 1\end{pmatrix} = \begin{pmatrix} 4\\ -3\\ -5\end{pmatrix}\]
  \textit{Gemeinsamer Punkt:} $x_3=1$ setzen.
  \[\begin{aligned} E:\ & 2x_1+x_2 = 3\\ F:\ & x_1-2x_2 = 4\end{aligned}\qquad\Longrightarrow\qquad 2\cdot\text{I}+\text{II}:\ 5x_1 = 10,\ x_1 = 2,\ x_2 = -1\]
  Kontrolle: $(2|-1|1)$ erfüllt $E$ ($4-1+1-4=0$) und $F$ ($2+2+2-6=0$).
  \[s:\ \vv{X} = \begin{pmatrix} 2\\ -1\\ 1\end{pmatrix}+t\cdot\begin{pmatrix} 4\\ -3\\ -5\end{pmatrix},\qquad t\in\mathds{R}\]
  (Alternativ: Eine Ebenengleichung nach einer Variablen auflösen und in die andere einsetzen.)
  \item Schnittwinkel von Ebenen gleich Winkel der Normalenvektoren (nicht stumpf):
  \[\cos{\varphi} = \dfrac{|\vv{n}_E\circ\vv{n}_F|}{|\vv{n}_E|\cdot|\vv{n}_F|} = \dfrac{|2-2+2|}{\sqrt{6}\cdot 3} = \dfrac{2}{3\sqrt{6}}\quad\Longrightarrow\quad\varphi\approx 74{,}21^{\circ}\]
\end{enumerate}
\end{loesung}

\Uebungsabschnitt{Abstände im Raum}

\begin{loesung}{2.9}{Abstände}
\begin{enumerate}[label=\alph*)]
  \item Hesse'sche Normalenform mit $|\vv{n}| = \sqrt{1+4+4} = 3$:
  \[d(P;E) = \dfrac{|3+2\cdot 4+2\cdot 5-6|}{3} = \dfrac{15}{3} = 5\]
  \item \textit{Verfahren:} Der Lotfußpunkt $F$ liegt auf $g$, und $\vv{QF}$ steht senkrecht auf dem Richtungsvektor $\vv{u}$ von $g$.
  \[\vv{QF} = \begin{pmatrix} 1+\lambda-3\\ 2-\lambda-0\\ 3+2\lambda-4\end{pmatrix} = \begin{pmatrix} \lambda-2\\ 2-\lambda\\ 2\lambda-1\end{pmatrix}\]
  \[\vv{QF}\circ\begin{pmatrix} 1\\ -1\\ 2\end{pmatrix} = (\lambda-2)-(2-\lambda)+2(2\lambda-1) = 6\lambda-6 = 0\quad\Longrightarrow\quad\lambda = 1\]
  Lotfußpunkt $F(2|1|5)$, $\vv{QF} = \begin{pmatrix} -1\\ 1\\ 1\end{pmatrix}$, also $d(Q;g) = |\vv{QF}| = \sqrt{3}\approx 1{,}73$.
  \item Merkekasten \glqq Abstand windschiefer Geraden\grqq{}: $d(g;h) = \left|(\vv{A}-\vv{B})\circ\vv{n_0}\right|$, wobei $\vv{n_0}$ ein Einheitsvektor senkrecht zu beiden Richtungsvektoren ist. (Dass $g$ und $h$ windschief sind, wurde in Ü~2.6\,b) gezeigt, denn $h=h_3$.)
  \[\vv{n} = \begin{pmatrix} 1\\ -1\\ 2\end{pmatrix}\times\begin{pmatrix} 1\\ 1\\ 0\end{pmatrix} = \begin{pmatrix} -1\cdot 0-2\cdot 1\\ 2\cdot 1-1\cdot 0\\ 1\cdot 1-(-1)\cdot 1\end{pmatrix} = \begin{pmatrix} -2\\ 2\\ 2\end{pmatrix},\qquad |\vv{n}| = \sqrt{12} = 2\sqrt{3}\]
  \[\vv{A}-\vv{B} = \begin{pmatrix} 1\\ 2\\ 3\end{pmatrix}-\begin{pmatrix} 0\\ 1\\ 0\end{pmatrix} = \begin{pmatrix} 1\\ 1\\ 3\end{pmatrix},\qquad d(g;h) = \dfrac{|-2+2+6|}{2\sqrt{3}} = \dfrac{6}{2\sqrt{3}} = \sqrt{3}\approx 1{,}73\]
\end{enumerate}
\end{loesung}

\Uebungsabschnitt{Anwendungen}

\begin{loesung}{2.10}{Kugel und Ebene}
\begin{enumerate}[label=\alph*)]
  \item Aus der Kugelgleichung $(x_1-m_1)^2+(x_2-m_2)^2+(x_3-m_3)^2 = R^2$ liest man ab: $M(2|1|-1)$, $R = \sqrt{25} = 5$.
  \begin{itemize}
    \item $P$: $(5-2)^2+(5-1)^2+(-1+1)^2 = 9+16+0 = 25 = R^2$, also liegt $P$ \textbf{auf} der Kugel.
    \item $Q$: $(4-2)^2+(3-1)^2+(0+1)^2 = 4+4+1 = 9 < 25$, also $|\vv{MQ}| = 3 < 5$, und $Q$ liegt \textbf{innerhalb} der Kugel.
  \end{itemize}
  \item Abstand Mittelpunkt--Ebene (Hesse'sche Normalenform, $|\vv{n}|=3$):
  \[d(M;E) = \dfrac{|2+2\cdot 1+2\cdot(-1)+7|}{3} = \dfrac{9}{3} = 3\]
  Wegen $d = 3 < R = 5$ schneidet $E$ die Kugel in einem Kreis.
  \item \textit{Mittelpunkt:} Lotfußpunkt von $M$ auf $E$. Die Lotgerade durch $M$ hat den Richtungsvektor $\vv{n}$, ihr allgemeiner Punkt ist $(2+\lambda|1+2\lambda|-1+2\lambda)$. Einsetzen in $E$:
  \[(2+\lambda)+2(1+2\lambda)+2(-1+2\lambda)+7 = 9+9\lambda = 0\quad\Longrightarrow\quad\lambda = -1\]
  Mittelpunkt des Schnittkreises: $M'(1|-1|-3)$.\\
  \textit{Radius} nach dem Satz des Pythagoras (Kugelradius, Abstand und Kreisradius bilden ein rechtwinkliges Dreieck):
  \[r = \sqrt{R^2-d^2} = \sqrt{25-9} = 4\]
\end{enumerate}
\end{loesung}

\begin{loesung}{2.11}{Spiegelung eines Punktes an einer Ebene}
\begin{enumerate}[label=\alph*)]
  \item Die Lotgerade verläuft durch $P$ in Richtung des Normalenvektors von $E$:
  \[\ell:\ \vv{X} = \begin{pmatrix} 3\\ 3\\ 3\end{pmatrix}+\lambda\cdot\begin{pmatrix} 1\\ 2\\ 2\end{pmatrix},\qquad\lambda\in\mathds{R}\]
  \item Schnitt von $\ell$ mit $E$:
  \[(3+\lambda)+2(3+2\lambda)+2(3+2\lambda)-6 = 9+9\lambda = 0\quad\Longrightarrow\quad\lambda = -1\]
  Lotfußpunkt $F(2|1|1)$. Abstand: $d(P;E) = |\vv{PF}| = \left|-1\cdot\begin{pmatrix} 1\\ 2\\ 2\end{pmatrix}\right| = 3$.\\
  (Kontrolle mit der Hesse'schen Normalenform: $\frac{|3+6+6-6|}{3} = 3$.)
  \item $F$ ist der Mittelpunkt von $\overline{PP'}$, also $\vv{P'} = \vv{P}+2\cdot\vv{PF}$. Das entspricht $\lambda = 2\cdot(-1) = -2$ in $\ell$:
  \[\vv{P'} = \begin{pmatrix} 3\\ 3\\ 3\end{pmatrix}-2\cdot\begin{pmatrix} 1\\ 2\\ 2\end{pmatrix} = \begin{pmatrix} 1\\ -1\\ -1\end{pmatrix},\qquad P'(1|-1|-1)\]
\end{enumerate}
\end{loesung}

\begin{loesung}{2.12}{Pyramide (Abiturniveau)}
\begin{enumerate}[label=\alph*)]
  \item $\vv{AB} = \begin{pmatrix} 4\\ 0\\ 0\end{pmatrix} = \vv{DC}$, also ist $ABCD$ ein Parallelogramm.\\
  $|\vv{AB}| = |\vv{AD}| = 4$ und $\vv{AB}\circ\vv{AD} = \begin{pmatrix} 4\\ 0\\ 0\end{pmatrix}\circ\begin{pmatrix} 0\\ 4\\ 0\end{pmatrix} = 0$. Ein Parallelogramm mit gleich langen Nachbarseiten und rechtem Winkel ist ein Quadrat.\\
  Die Grundfläche liegt in der $x_1x_2$-Ebene, die Höhe ist also die $x_3$-Koordinate von $S$: $h=6$.
  \[V = \dfrac{1}{3}\cdot G\cdot h = \dfrac{1}{3}\cdot 16\cdot 6 = 32\ \text{VE}\]
  \item Spannvektoren $\vv{BC} = \begin{pmatrix} 0\\ 4\\ 0\end{pmatrix}$ und $\vv{BS} = \begin{pmatrix} -2\\ 2\\ 6\end{pmatrix}$.
  \[\vv{BC}\times\vv{BS} = \begin{pmatrix} 4\cdot 6-0\cdot 2\\ 0\cdot(-2)-0\cdot 6\\ 0\cdot 2-4\cdot(-2)\end{pmatrix} = \begin{pmatrix} 24\\ 0\\ 8\end{pmatrix} = 8\cdot\begin{pmatrix} 3\\ 0\\ 1\end{pmatrix}\]
  Ansatz $3x_1+x_3-c=0$, Punktprobe mit $B(4|0|0)$: $c=12$, also $L: 3x_1+x_3-12=0$.\\
  Kontrolle: $C$: $12+0-12=0$, $S$: $6+6-12=0$.
  \item Die Grundfläche liegt in der $x_1x_2$-Ebene mit Normalenvektor $\vv{e}_3 = \begin{pmatrix} 0\\ 0\\ 1\end{pmatrix}$. Schnittwinkel von Ebenen:
  \[\cos{\varphi} = \dfrac{\left|\begin{pmatrix} 3\\ 0\\ 1\end{pmatrix}\circ\begin{pmatrix} 0\\ 0\\ 1\end{pmatrix}\right|}{\sqrt{10}\cdot 1} = \dfrac{1}{\sqrt{10}}\quad\Longrightarrow\quad\varphi\approx 71{,}57^{\circ}\]
  \item Hesse'sche Normalenform mit $|\vv{n}_L| = \sqrt{10}$:
  \[d(A;L) = \dfrac{|0+0-12|}{\sqrt{10}} = \dfrac{12}{\sqrt{10}} = \dfrac{6\sqrt{10}}{5}\approx 3{,}79\]
  \item Mittelpunkt: $\vv{N} = \frac{1}{2}(\vv{B}+\vv{C})$, also $N(4|2|0)$.
  \[\vv{NS} = \begin{pmatrix} -2\\ 0\\ 6\end{pmatrix},\qquad |\vv{NS}| = \sqrt{4+36} = \sqrt{40} = 2\sqrt{10}\approx 6{,}32\]
  Wegen $\vv{NS}\circ\vv{BC} = 0$ steht $\overline{NS}$ senkrecht auf $\overline{BC}$, ist also die Höhe des Dreiecks $BCS$ zur Grundseite $\overline{BC}$.
  \[A_{BCS} = \dfrac{1}{2}\cdot|\vv{BC}|\cdot|\vv{NS}| = \dfrac{1}{2}\cdot 4\cdot 2\sqrt{10} = 4\sqrt{10}\approx 12{,}65\ \text{FE}\]
  \textit{Bestätigung mit dem Vektorprodukt} (aus b)):
  \[A_{BCS} = \dfrac{1}{2}\left|\vv{BC}\times\vv{BS}\right| = \dfrac{1}{2}\sqrt{24^2+0^2+8^2} = \dfrac{1}{2}\sqrt{640} = \dfrac{1}{2}\cdot 8\sqrt{10} = 4\sqrt{10}\]
\end{enumerate}
\end{loesung}

\Uebungsabschnitt{Abiturniveau \textendash\ Teil A (ohne Hilfsmittel)}

\begin{loesung}{2.13}{Ebene und Abstand}
\begin{enumerate}[label=\alph*)]
  \item Je zwei Koordinaten null setzen (Spurpunkte mit den Achsen):
  \[x_1\text{-Achse: } 2x_1 = 6\Rightarrow S_1(3|0|0),\quad x_2\text{-Achse: } -x_2 = 6\Rightarrow S_2(0|-6|0),\quad x_3\text{-Achse: } 2x_3 = 6\Rightarrow S_3(0|0|3)\]
  \item Hesse'sche Normalenform mit $|\vv{n}| = \sqrt{4+1+4} = 3$:
  \[d(O;E) = \dfrac{|0-0+0-6|}{3} = 2\]
  \item Parallele Ebenen haben denselben Normalenvektor: $E_c: 2x_1-x_2+2x_3-c = 0$. Der Abstand zweier paralleler Ebenen ist der Abstand eines Punktes der einen Ebene von der anderen, z.\,B. von $S_1(3|0|0)\in E$:
  \[\dfrac{|6-c|}{3} = 2\quad\Longleftrightarrow\quad|6-c| = 6\quad\Longleftrightarrow\quad c = 0\ \text{oder}\ c = 12\]
  Gesuchte Ebenen: $2x_1-x_2+2x_3 = 0$ (durch den Ursprung, passend zu b)) und $2x_1-x_2+2x_3-12 = 0$.
\end{enumerate}
\end{loesung}

\begin{loesung}{2.14}{Gerade und Kugel}
\begin{enumerate}[label=\alph*)]
  \item Der Aufpunkt von $g$ ist $M$ (Parameterwert $\lambda=0$), also verläuft $g$ durch $M$.\\
  Der Richtungsvektor hat die Länge $\sqrt{4+1+4} = 3 = r$. Die Schnittpunkte haben daher von $M$ den Abstand $r$, wenn $|\lambda|\cdot 3 = 3$, also $\lambda = \pm 1$:
  \[\lambda = 1:\ P(3|3|1),\qquad\lambda = -1:\ Q(-1|1|-3)\]
  \item Die Tangentialebene steht senkrecht auf dem Berührradius $\vv{MP} = \begin{pmatrix} 2\\ 1\\ 2\end{pmatrix}$, dieser ist also Normalenvektor. Ansatz $2x_1+x_2+2x_3-c = 0$, Punktprobe mit $P$: $c = 6+3+2 = 11$.
  \[T: 2x_1+x_2+2x_3-11 = 0\]
\end{enumerate}
\end{loesung}

\begin{loesung}{2.15}{Gerade und Ebene mit Parameter}
\begin{enumerate}[label=\alph*)]
  \item Normalenvektor $\vv{n} = \begin{pmatrix} 1\\ 2\\ -2\end{pmatrix}$, Richtungsvektor $\vv{u} = \begin{pmatrix} 2\\ -1\\ 0\end{pmatrix}$:
  \[\vv{n}\circ\vv{u} = 2-2+0 = 0\]
  Der Richtungsvektor ist orthogonal zum Normalenvektor, also verläuft jede Gerade $g_a$ parallel zu $E$ (oder liegt in $E$).
  \item Eine parallele Gerade liegt genau dann in $E$, wenn ihr Aufpunkt $(0|0|a)$ in $E$ liegt (Punktprobe):
  \[0+0-2a-4 = 0\quad\Longleftrightarrow\quad a = -2\]
\end{enumerate}
\end{loesung}

\begin{loesung}{2.16}{Rechtwinkliges Dreieck}
\begin{enumerate}[label=\alph*)]
  \item \[\vv{AB} = \begin{pmatrix} 2\\ 2\\ -1\end{pmatrix},\quad\vv{AC} = \begin{pmatrix} 2\\ -1\\ 2\end{pmatrix},\quad\vv{AB}\circ\vv{AC} = 4-2-2 = 0\]
  Also ist der Winkel bei $A$ ein rechter. Längen: $|\vv{AB}| = \sqrt{4+4+1} = 3$ und $|\vv{AC}| = \sqrt{4+1+4} = 3$, also gleichschenklig.
  \item Rechtwinkliges Dreieck: $A = \dfrac{1}{2}\cdot|\vv{AB}|\cdot|\vv{AC}| = \dfrac{1}{2}\cdot 3\cdot 3 = 4{,}5$.
  \item Nach a) bilden $\vv{AB}$ und $\vv{AC}$ gleich lange, orthogonale Seiten. Im Quadrat $ABDC$ gilt $\vv{BD} = \vv{AC}$:
  \[\vv{D} = \vv{B}+\vv{AC} = \begin{pmatrix} 3\\ 2\\ 1\end{pmatrix}+\begin{pmatrix} 2\\ -1\\ 2\end{pmatrix} = \begin{pmatrix} 5\\ 1\\ 3\end{pmatrix}\quad\Longrightarrow\quad D(5|1|3)\]
  Kontrolle: $\vv{C}+\vv{AB} = (3+2|-1+2|4-1) = (5|1|3)$. \checkmark
\end{enumerate}
\end{loesung}

\Uebungsabschnitt{Abiturniveau \textendash\ Teil B (mit Hilfsmitteln)}

\begin{loesung}{2.17}{Dachfläche mit Antenne}
\begin{enumerate}[label=\alph*)]
  \item $\vv{PQ} = \begin{pmatrix} 10\\ 0\\ 0\end{pmatrix}$, $\vv{PR} = \begin{pmatrix} 0\\ 8\\ 6\end{pmatrix}$. Normalenvektor über das Vektorprodukt:
  \[\vv{PQ}\times\vv{PR} = \begin{pmatrix} 0\cdot 6-0\cdot 8\\ 0\cdot 0-10\cdot 6\\ 10\cdot 8-0\cdot 0\end{pmatrix} = \begin{pmatrix} 0\\ -60\\ 80\end{pmatrix} = 20\cdot\begin{pmatrix} 0\\ -3\\ 4\end{pmatrix}\]
  Ansatz $-3x_2+4x_3-c = 0$, Punktprobe mit $P$: $c = 12$. Also $D: -3x_2+4x_3-12 = 0$.
  \item Winkel zwischen $D$ und dem Boden ($x_1x_2$-Ebene, Normalenvektor $\vv{e}_3$):
  \[\cos{\varphi} = \dfrac{|4|}{\sqrt{9+16}\cdot 1} = \dfrac{4}{5}\quad\Longrightarrow\quad\varphi\approx 36{,}87^{\circ}\]
  \item $F$: $-12+24-12 = 0$, also $F\in D$. Die Spitze ist $S(5|4|8)$. Hesse'sche Normalenform ($|\vv{n}| = 5$):
  \[d(S;D) = \dfrac{|-12+32-12|}{5} = \dfrac{8}{5} = 1{,}6\ (\text{m})\]
  Der Abstand ist kleiner als die Antennenlänge, weil die Antenne nicht senkrecht auf der schrägen Dachfläche steht.
  \item Lichtstrahl durch $S$: $\vv{X} = \begin{pmatrix} 5\\ 4\\ 8\end{pmatrix}+\lambda\cdot\begin{pmatrix} 1\\ 0\\ -2\end{pmatrix}$. Einsetzen in $D$:
  \[-3\cdot 4+4\cdot(8-2\lambda)-12 = 8-8\lambda = 0\quad\Longrightarrow\quad\lambda = 1\]
  Schattenpunkt $S'(6|4|6)$. Der Schatten der Antenne ist die Strecke $\overline{FS'}$ mit $\vv{FS'} = \begin{pmatrix} 1\\ 0\\ 0\end{pmatrix}$, also $1\,\text{m}$ lang.
\end{enumerate}
\end{loesung}

\begin{loesung}{2.18}{Flugbahnen}
\begin{enumerate}[label=\alph*)]
  \item Pro Minute legen die Flugzeuge die Länge des Richtungsvektors zurück:
  \[\left|\begin{pmatrix} 4\\ 3\\ 0\end{pmatrix}\right| = \left|\begin{pmatrix} -4\\ 3\\ 0\end{pmatrix}\right| = 5\ \tfrac{\text{km}}{\text{min}} = 300\ \tfrac{\text{km}}{\text{h}}\]
  \item Die Richtungsvektoren sind nicht parallel ($4:(-4)\neq 3:3$). Gleichsetzen der Bahngeraden (Parameter $r$, $s$):
  \[\text{I: } 4r = 24-4s,\qquad\text{II: } 3r = 3s,\qquad\text{III: } 8 = 7\]
  Gleichung III ist ein Widerspruch, es gibt keinen Schnittpunkt. Die Bahnen sind weder parallel noch schneiden sie sich, sie sind also windschief.
  \item Abstand windschiefer Geraden:
  \[\vv{n} = \begin{pmatrix} 4\\ 3\\ 0\end{pmatrix}\times\begin{pmatrix} -4\\ 3\\ 0\end{pmatrix} = \begin{pmatrix} 0\\ 0\\ 24\end{pmatrix},\quad\vv{n_0} = \begin{pmatrix} 0\\ 0\\ 1\end{pmatrix},\quad d = \left|\left(\begin{pmatrix} 0\\ 0\\ 8\end{pmatrix}-\begin{pmatrix} 24\\ 0\\ 7\end{pmatrix}\right)\circ\vv{n_0}\right| = 1\]
  Die Bahnen haben den Abstand $1\,\text{km}$ (beide verlaufen waagrecht, in $8\,\text{km}$ bzw. $7\,\text{km}$ Höhe).
  \item Verbindungsvektor der Positionen zum gleichen Zeitpunkt:
  \[\vv{X_1}(t)-\vv{X_2}(t) = \begin{pmatrix} 8t-24\\ 0\\ 1\end{pmatrix},\qquad d(t) = \sqrt{(8t-24)^2+1}\]
  $d(t)$ ist minimal, wenn $(8t-24)^2$ minimal ist, also für $t=3$, mit $d(3) = 1$. Nach $3$ Minuten befinden sich beide Flugzeuge über demselben Punkt $(12|9)$ und kommen sich bis auf $1\,\text{km}$ nahe. Der kleinste Abstand der Flugzeuge stimmt hier mit dem Abstand der Bahnen überein, weil sie die Stellen des kürzesten Abstands gleichzeitig erreichen. (Im Allgemeinen ist der kleinste Abstand der Flugzeuge größer als der Abstand der Bahnen.)
  \item Winkel zwischen den Richtungsvektoren (nicht stumpf):
  \[\cos{\varphi} = \dfrac{|-16+9+0|}{5\cdot 5} = \dfrac{7}{25}\quad\Longrightarrow\quad\varphi\approx 73{,}74^{\circ}\]
\end{enumerate}
\end{loesung}

\begin{loesung}{2.19}{Schnitt durch einen Würfel}
\begin{enumerate}[label=\alph*)]
  \item $B$, $D$, $E$ liegen auf den Koordinatenachsen, jeweils im Abstand $6$ vom Ursprung (Achsenabschnittsform). Ansatz $x_1+x_2+x_3-c=0$; Probe mit $B$: $c=6$, mit $D$ und $E$ ebenfalls erfüllt.
  \[L: x_1+x_2+x_3-6 = 0\]
  \item Hesse'sche Normalenform mit $|\vv{n}| = \sqrt{3}$:
  \[d(A;L) = \dfrac{|-6|}{\sqrt{3}} = 2\sqrt{3}\approx 3{,}46,\qquad d(G;L) = \dfrac{|18-6|}{\sqrt{3}} = 4\sqrt{3}\approx 6{,}93\]
  \item $\vv{AG} = \begin{pmatrix} 6\\ 6\\ 6\end{pmatrix} = 6\cdot\vv{n}$ ist ein Vielfaches des Normalenvektors, also steht $\overline{AG}$ senkrecht auf $L$.\\
  Schnittpunkt: $\vv{X} = \lambda\cdot\begin{pmatrix} 6\\ 6\\ 6\end{pmatrix}$ in $L$: $18\lambda = 6$, $\lambda = \frac{1}{3}$, also $T(2|2|2)$. $L$ teilt $\overline{AG}$ im Verhältnis $1:2$ (von $A$ aus). Das passt zu b): $2\sqrt{3}:4\sqrt{3} = 1:2$.
  \item Die Pyramide $ABDE$ wird von $\vv{AB}$, $\vv{AD}$, $\vv{AE}$ aufgespannt (paarweise senkrecht, je Länge $6$):
  \[V_{\text{Pyramide}} = \dfrac{1}{6}\left|(\vv{AB}\times\vv{AD})\circ\vv{AE}\right| = \dfrac{1}{6}\cdot 6^3 = 36,\qquad V_{\text{Würfel}} = 216\]
  Anteil: $\frac{36}{216} = \frac{1}{6}$.
\end{enumerate}
\end{loesung}

\begin{loesung}{2.20}{Kugel auf einer schiefen Ebene}
\begin{enumerate}[label=\alph*)]
  \item Hesse'sche Normalenform ($|\vv{n}| = 3$):
  \[d(M;E) = \dfrac{|8-1+8-6|}{3} = \dfrac{9}{3} = 3 = r\]
  Der Abstand des Mittelpunkts von $E$ ist gleich dem Radius, also berührt $E$ die Kugel.\\
  Berührpunkt: Lotfußpunkt von $M$. Die Lotgerade $\vv{X} = \vv{M}+\lambda\vv{n}$ eingesetzt in $E$:
  \[2(4+2\lambda)-(1-\lambda)+2(4+2\lambda)-6 = 9+9\lambda = 0\quad\Longrightarrow\quad\lambda = -1,\qquad B(2|2|2)\]
  \item Der zweite Berührpunkt liegt auf der Lotgeraden auf der anderen Seite von $M$: $\lambda = +1$, also $B'(6|0|6)$. Ebene mit gleichem Normalenvektor durch $B'$:
  \[2x_1-x_2+2x_3-c = 0,\quad c = 12-0+12 = 24,\qquad E': 2x_1-x_2+2x_3-24 = 0\]
  \item Der höchste Punkt liegt senkrecht über $M$ im Abstand $r$: $\vv{M}+3\cdot\vv{e}_3$, also $(4|1|7)$.
\end{enumerate}
\end{loesung}

\begin{loesung}{2.21}{Zelt}
\begin{enumerate}[label=\alph*)]
  \item $W_1$: $\vv{AD} = \begin{pmatrix} 0\\ 6\\ 0\end{pmatrix}$, $\vv{AF} = \begin{pmatrix} 2\\ 0\\ 2\end{pmatrix}$, Vektorprodukt $\vv{AD}\times\vv{AF} = \begin{pmatrix} 12\\ 0\\ -12\end{pmatrix}$, also $\vv{n_1} = \begin{pmatrix} 1\\ 0\\ -1\end{pmatrix}$. Mit $A$: $W_1: x_1-x_3 = 0$.\\
  $W_2$: $\vv{BC} = \begin{pmatrix} 0\\ 6\\ 0\end{pmatrix}$, $\vv{BF} = \begin{pmatrix} -2\\ 0\\ 2\end{pmatrix}$, $\vv{BC}\times\vv{BF} = \begin{pmatrix} 12\\ 0\\ 12\end{pmatrix}$, also $\vv{n_2} = \begin{pmatrix} 1\\ 0\\ 1\end{pmatrix}$. Mit $B$: $W_2: x_1+x_3-4 = 0$.\\
  Probe: $F(2|0|2)$ erfüllt $2-2=0$ und $2+2-4=0$. \checkmark
  \item Schnittwinkel zweier Ebenen über die Normalenvektoren: $\vv{n_1}\circ\vv{n_2} = 1+0-1 = 0$, also $\cos{\varphi} = 0$ und $\varphi = 90^\circ$. Die Seitenwände stehen senkrecht aufeinander.
  \item Das Zelt ist ein liegendes Prisma mit dem Dreieck $ABF$ als Grundfläche (Grundseite $4$, Höhe $2$) und der Länge $6$:
  \[V = \dfrac{1}{2}\cdot 4\cdot 2\cdot 6 = 24\]
  Das Zelt hat ein Volumen von $24\,\text{m}^3$.
  \item Die Stange liegt auf der Geraden $x_1 = 1$, $x_2 = 3$; wegen $x_1 = 1<2$ trifft sie die Wand $W_1$. Einsetzen: $1-x_3 = 0\Longleftrightarrow x_3 = 1$. Die Stange berührt die Wand in $1\,\text{m}$ Höhe im Punkt $(1|3|1)$.
  \item Hessesche Normalform von $W_1$: $\dfrac{x_1-x_3}{\sqrt{2}} = 0$.
  \[d(L,W_1) = \dfrac{|2-1|}{\sqrt{2}} = \dfrac{\sqrt{2}}{2}\approx 0{,}71\]
  Die Lampe ist etwa $0{,}71\,\text{m}$ von der Seitenwand entfernt.
\end{enumerate}
\end{loesung}

\begin{loesung}{2.22}{Reflektierter Laserstrahl}
\begin{enumerate}[label=\alph*)]
  \item Strahl $s: \vv{X} = \begin{pmatrix} 0\\ 0\\ 6\end{pmatrix}+\lambda\cdot\begin{pmatrix} 1\\ 1\\ -2\end{pmatrix}$. Boden: $x_3 = 0\Longleftrightarrow 6-2\lambda = 0\Longleftrightarrow\lambda = 3$. Also $R(3|3|0)$.
  \item Bei der Reflexion an der $x_1x_2$-Ebene bleiben die zur Ebene parallelen Komponenten erhalten, die Komponente in Normalenrichtung ($x_3$) wechselt das Vorzeichen: aus $\begin{pmatrix} 1\\ 1\\ -2\end{pmatrix}$ wird $\begin{pmatrix} 1\\ 1\\ 2\end{pmatrix}$.\\
  Reflektierter Strahl $r: \vv{X} = \begin{pmatrix} 3\\ 3\\ 0\end{pmatrix}+\mu\cdot\begin{pmatrix} 1\\ 1\\ 2\end{pmatrix}$. Wand: $3+\mu = 6\Longleftrightarrow\mu = 3$, also $W(6|6|6)$.
  \item Winkel zwischen Gerade und Ebene mit $\vv{n} = \begin{pmatrix} 0\\ 0\\ 1\end{pmatrix}$:
  \[\sin{\alpha} = \dfrac{|\vv{u}\circ\vv{n}|}{|\vv{u}|\cdot|\vv{n}|} = \dfrac{2}{\sqrt{6}}\quad\Longrightarrow\quad\alpha\approx 54{,}7^\circ\]
\end{enumerate}
\end{loesung}

\LoesungskapitelEnde
